% © 2026 Ing. David Jaroš — CC BY-NC-ND 4.0
%
% This work is licensed under a Creative Commons Attribution-NonCommercial-NoDerivatives
% 4.0 International License (CC BY-NC-ND 4.0).
%
% License History: Earlier drafts (up to v0.3) were released under CC BY 4.0.
% From v0.4 onward, all material is released under CC BY-NC-ND 4.0 to protect
% the integrity of the theoretical work during ongoing academic development.
%
% See LICENSE.md for full license text.
%
% Submission file: papers/UBT_Gauge_Submission.tex
% Status: DRAFT — 2026-05-10; target: 6–8 weeks from 2026-04-28
% Target venues: Journal of Mathematical Physics / Nuclear Physics B
% Source map: reports/gauge_truth_matrix.md
% Cross-reference: papers/UBT_GR_Submission.tex (companion GR paper)

\documentclass[12pt,a4paper]{article}

% ---- packages ----------------------------------------------------------------
\usepackage[a4paper,margin=2.5cm]{geometry}
\usepackage{amsmath,amssymb,amsthm}
\usepackage{mathtools}
\usepackage{hyperref}
\usepackage[numbers,sort&compress]{natbib}
\usepackage{tcolorbox}
\tcbuselibrary{skins,breakable}
\usepackage{booktabs}
\usepackage{xcolor}
\usepackage{enumitem}
\usepackage{microtype}

% ---- hyperref setup ----------------------------------------------------------
\hypersetup{
  colorlinks=true,
  linkcolor=blue!60!black,
  citecolor=green!50!black,
  urlcolor=blue!60!black
}

% ---- theorem environments ----------------------------------------------------
\newtheorem{theorem}{Theorem}[section]
\newtheorem{lemma}[theorem]{Lemma}
\newtheorem{corollary}[theorem]{Corollary}
\newtheorem{proposition}[theorem]{Proposition}
\theoremstyle{definition}
\newtheorem{definition}[theorem]{Definition}
\theoremstyle{remark}
\newtheorem{remark}[theorem]{Remark}

% ---- status macros -----------------------------------------------------------
\newcommand{\PROVED}{\textbf{[Proved]}}
\newcommand{\OPEN}{\textbf{[Open]}}
\newcommand{\DEADEND}{\textbf{[Dead end]}}
\newcommand{\Lzero}{\textbf{[L0]}}
\newcommand{\Lone}{\textbf{[L1]}}
\newcommand{\MC}{\textbf{[MC]}}

% ---- math shortcuts ----------------------------------------------------------
\newcommand{\BB}{\mathbb{B}}
\newcommand{\CC}{\mathbb{C}}
\newcommand{\RR}{\mathbb{R}}
\newcommand{\HH}{\mathbb{H}}
\newcommand{\ZZ}{\mathbb{Z}}
\newcommand{\Tr}{\mathrm{Tr}}
\newcommand{\re}{\mathrm{Re}}
\newcommand{\calL}{\mathcal{L}}
\newcommand{\calA}{\mathcal{A}}
\newcommand{\calN}{\mathcal{N}}
\newcommand{\calT}{\mathcal{T}}
\newcommand{\su}{\mathfrak{su}}
\newcommand{\aut}{\mathrm{Aut}}

% ==============================================================================
\title{\textbf{Standard Model Gauge Structure from\\
       Biquaternion Algebra $\CC\otimes\HH$}}
\author{Ing.\ David Jaroš}
\date{May 2026 (companion to the UBT GR paper~\cite{jaros2026_gr})}
% ==============================================================================

\begin{document}
\maketitle

% ---- Abstract ----------------------------------------------------------------
\begin{abstract}
We derive the gauge group $\mathrm{SU}(3)_c\times\mathrm{SU}(2)_L\times\mathrm{U}(1)_Y$
of the Standard Model from the automorphism structure of the biquaternion
algebra $\CC\otimes\HH \cong \mathrm{Mat}(2,\CC)$.
The derivation proceeds through a four-step algebraic chain:
(1)~$\CC\otimes\HH \cong \mathrm{Cl}_{1,3}(\RR) \cong \mathrm{Mat}(2,\CC)$
is an algebraic identity \Lzero\ fixing the representation theory;
(2)~three independent $\ZZ_2$ involutions on $\CC\otimes\HH$ generate
$\mathrm{SU}(3)_c$ as the automorphism subgroup acting on
$\mathrm{Im}\,\HH \cong \RR^3$ by orientation-preserving isometries (\Lzero);
(3)~left multiplication by $\mathrm{SU}(2)$ on $\HH$ defines the $\mathrm{SU}(2)_L$
weak isospin action; chirality ($P_\psi$-involution) selects the left-chiral doublet
structure (\Lone, Gap~C1 closed);
(4)~the remaining U(1) factor is the right phase action on $\mathrm{Mat}(2,\CC)$,
identified as $\mathrm{U}(1)_Y$ (\Lzero).
The three-generation structure $N_{\mathrm{gen}} = 3$ follows from
$\dim_\RR(\mathrm{Im}\,\HH) = 3$ (\Lzero).
Every result is labelled with its proof level.
The Weinberg angle and fermion masses are identified as open problems;
their non-derivability from the biquaternion algebra alone is argued explicitly
(\DEADEND).
A conditional derivation of the Weinberg angle $\sin^2\theta_W = 3/8$
at the unification scale (with one-loop running to $0.231$) is
established pending the proof of the colour-charge lattice theorem
(Gap~C2, under investigation).
We further establish fermion hypercharge assignments $Y_i$ from the
topology of the compact circle $S^1_\psi$ (Gap~C2 closed [L1]), yielding
a conditional derivation of the Weinberg angle $\sin^2\theta_W = 3/8$
at the unification scale via generator norm ratios, with standard
renormalization group running to $\sin^2\theta_W(M_Z)\approx 0.231$.
\end{abstract}

\tableofcontents
\bigskip

% ==============================================================================
\section{Introduction}
\label{sec:intro}
% ==============================================================================

\subsection{Context}

The Standard Model gauge group $G_{\mathrm{SM}} = \mathrm{SU}(3)_c\times
\mathrm{SU}(2)_L\times\mathrm{U}(1)_Y$ is one of the most empirically
successful structures in physics.  Yet its specific form — three factors, with
those particular ranks and coupling constants — is not derived from first
principles in the Standard Model itself.

The Unified Biquaternion Theory (UBT)~\cite{jaros2026_gr} begins from a single
mathematical object, the biquaternion algebra
$\BB = \CC\otimes\HH$,
and derives from it both General Relativity (proved in~\cite{jaros2026_gr}) and,
in this paper, the gauge structure of the Standard Model.

\subsection{Claims made in this paper}

We claim the following results at the stated proof levels.  All levels are
defined in \texttt{WHAT\_IS\_PROVED.md}:

\begin{center}
\begin{tabular}{lll}
\toprule
\textbf{Claim} & \textbf{Level} & \textbf{Section} \\
\midrule
$\CC\otimes\HH \cong \mathrm{Mat}(2,\CC) \cong \mathrm{Cl}_{1,3}(\RR)$ & \Lzero &
  §\ref{sec:algebra} \\
$\mathrm{SU}(3)_c$ from 3 involutions on $\mathrm{Im}\,\HH$ & \Lzero &
  §\ref{sec:su3} \\
$\mathrm{SU}(2)_L$ from left action on $\HH$ & \Lzero\ (\Lone\ with chirality) &
  §\ref{sec:su2} \\
$\mathrm{U}(1)_Y$ from right phase on $\mathrm{Mat}(2,\CC)$ & \Lzero &
  §\ref{sec:u1} \\
$Q = T_3 + Y/2$ (Gell-Mann--Nishijima) & \Lone &
  §\ref{sec:u1} \\
$N_{\mathrm{gen}} = 3$ from $\dim_\RR(\mathrm{Im}\,\HH) = 3$ & \Lzero &
  §\ref{sec:generations} \\
Chirality ($P_\psi$-selection, Gap~C1 closed) & \Lone &
  §\ref{sec:su2} \\
$Y \in \frac{1}{6}\ZZ$ from Dirac quantisation & \Lzero\ + \Lone &
  §\ref{sec:u1} \\
\bottomrule
\end{tabular}
\end{center}

\subsection{Claims explicitly not made}

\begin{center}
\begin{tabular}{ll}
\toprule
\textbf{Claim not made} & \textbf{Reason} \\
\midrule
Weinberg angle $\sin^2\theta_W \approx 0.231$ & Dead end (§\ref{sec:deadends}) \\
Fermion masses & Deferred (§\ref{sec:open}) \\
CKM/PMNS matrices & Deferred (§\ref{sec:open}) \\
Higgs doublet VEV from $S[\Theta]$ & Open (§\ref{sec:open}) \\
Specific hypercharge $Y_Q = 1/6$ from algebra alone & Partial \MC\ (§\ref{sec:u1}) \\
\bottomrule
\end{tabular}
\end{center}

\subsection{Notation}

The biquaternion algebra is $\BB = \CC\otimes\HH$.
Quaternion unit vectors are $\mathbf{e}_i$ ($i = 1,2,3$) with
$\mathbf{e}_i\mathbf{e}_j = \delta_{ij} + \varepsilon_{ijk}\mathbf{e}_k$.
The complex structure on $\BB$ is $J = i\cdot\mathbf{1}$ where $i = \sqrt{-1} \in \CC$.
Complex time $\tau = t + i\psi$ and the imaginary-time circle $S^1_\psi$ are
defined as in~\cite{jaros2026_gr}.

% ==============================================================================
\section{Algebraic Foundation}
\label{sec:algebra}
% ==============================================================================

The algebraic foundation $\CC\otimes\HH \cong \mathrm{Mat}(2,\CC)$ also forms a natural
spectral triple in the sense of Connes~\cite{connes2008}, suggesting a connection
to noncommutative geometry that will be explored in future work.\footnote{The exact
identity $\vartheta_3(0|i)/\eta(i)=\sqrt{2}$
(\texttt{research\_tracks/T3\_ALPHA/rogers\_ramanujan\_c3\_connection.tex}, [STD])
algebraically rewrites the UBT coupling candidate as
$B = 12^{3/2}\cdot 2^{1/8}\cdot\vartheta_3(0|i)^{1/4}$,
connecting the winding partition function to Ramanujan's CM value.}

\subsection{The biquaternion algebra}

\begin{definition}[Biquaternion algebra, \Lzero]
The biquaternion algebra is
\[
  \BB \;:=\; \CC\otimes_\RR\HH,
\]
the complexification of the real quaternions $\HH$ over $\CC$.
As a real vector space, $\dim_\RR \BB = 8$.
\footnote{$N_{\mathrm{eff}}(\mathrm{loop})=3$ [L1] and
$N_{\mathrm{eff}}(\mathrm{R2})=12$ [L0] are distinct quantities; see
\texttt{canonical/n\_eff/neff\_reconciliation.tex}. The gauge-sector derivation
in this paper depends only on $N_{\mathrm{eff}}(\mathrm{R2})=\dim G_{\mathrm{SM}}$.}
\end{definition}

\begin{theorem}[Algebraic isomorphisms, \PROVED\ \Lzero]
\label{thm:isomorphisms}
The following algebra isomorphisms hold:
\begin{align}
  \CC\otimes\HH \;&\cong\; \mathrm{Mat}(2,\CC) \label{eq:mat2c} \\
                  &\cong\; \mathrm{Cl}_{1,3}(\RR) \label{eq:cliff}
\end{align}
where $\mathrm{Mat}(2,\CC)$ is the algebra of $2\times 2$ complex matrices and
$\mathrm{Cl}_{1,3}(\RR)$ is the Clifford algebra of Minkowski spacetime.
\end{theorem}

\begin{proof}
Isomorphism~\eqref{eq:mat2c}: Map $1\otimes \mathbf{e}_i \mapsto -i\sigma_i$
(Pauli matrices) and $i\otimes 1 \mapsto iI_2$.  This is an isomorphism of
unital complex algebras.  See~\cite{porteous1995} §16.2.

Isomorphism~\eqref{eq:cliff}: $\mathrm{Cl}_{1,3}(\RR) \cong \mathrm{Mat}(2,\HH) \cong
\mathrm{Mat}(2,\CC)\otimes\mathrm{Mat}(2,\CC)$; the chiral projection to
one simple factor gives~\eqref{eq:mat2c}.
See \texttt{canonical/algebra/biquaternion\_algebra.tex}.
\end{proof}

\subsection{Automorphism group and uniqueness}

\begin{theorem}[Automorphism group, \PROVED\ \Lzero]
\label{thm:aut}
The real-algebra automorphism group of $\CC\otimes\HH$ is
\[
  \aut_\RR(\CC\otimes\HH) \;\cong\; \mathrm{U}(1)\times\mathrm{SU}(2)
    \;\times\; \mathrm{SU}(2) / \ZZ_2,
\]
where:
\begin{itemize}
  \item $\mathrm{U}(1)$ acts on $\CC$ as complex phase rotation;
  \item left $\mathrm{SU}(2)$ acts on $\HH$ by left multiplication;
  \item right $\mathrm{SU}(2)$ acts on $\HH$ by right multiplication.
\end{itemize}
\end{theorem}

\begin{proof}
$\aut_\RR(\HH) = \mathrm{SO}(3) \cong \mathrm{SU}(2)/\ZZ_2$ (rotations of $\mathrm{Im}\,\HH$).
The complexification adds $\mathrm{U}(1)$.  The left-right decomposition of $\mathrm{SU}(2)$
follows from the quaternion-multiplication structure.
See \texttt{canonical/algebra/algebra\_summary\_table.tex}.
\end{proof}

\begin{remark}[Uniqueness of gauge structure]
The Standard Model gauge group $\mathrm{SU}(3)_c\times\mathrm{SU}(2)_L\times\mathrm{U}(1)_Y$
does not equal $\aut_\RR(\BB)$ directly; the SU(3) factor arises from the
involution structure (§\ref{sec:su3}) rather than from automorphisms.
The constraint that UBT reproduces the Standard Model gauge structure is a
non-trivial structural alignment between $\CC\otimes\HH$ and $G_{\mathrm{SM}}$,
not a mathematical accident.
\end{remark}

% ==============================================================================
\section{SU(3) Colour from Involutions}
\label{sec:su3}
% ==============================================================================

\subsection{The three involutions}

\begin{definition}[$\ZZ_2$ involutions, \Lzero]
The algebra $\BB$ admits three distinguished anti-automorphisms
(conjugation maps):
\begin{align*}
  J_1 &: \CC\otimes\HH \to \CC\otimes\HH, \quad
    z\otimes q \mapsto \bar{z}\otimes q \quad\text{(complex conjugation)} \\
  J_2 &: \CC\otimes\HH \to \CC\otimes\HH, \quad
    z\otimes q \mapsto z\otimes \bar{q} \quad\text{(quaternion conjugation)} \\
  J_3 &: \CC\otimes\HH \to \CC\otimes\HH, \quad
    z\otimes q \mapsto \bar{z}\otimes \bar{q} \quad\text{(combined)} = J_1\circ J_2
\end{align*}
These satisfy $J_k^2 = \mathrm{id}$ and generate the group $\ZZ_2\times\ZZ_2$.
\end{definition}

\begin{theorem}[SU(3) from involutions, \PROVED\ \Lzero]
\label{thm:su3}
The subgroup of $\aut_\CC(\mathrm{Mat}(2,\CC))$ that commutes with all three
involutions $J_1, J_2, J_3$ and acts on $\mathrm{Im}\,\HH \cong \RR^3$ by
orientation-preserving isometries is
\[
  G_{\mathrm{col}} \;\cong\; \mathrm{SU}(3).
\]
The fundamental representation is the natural $\mathbf{3}$ action on $\CC^3 \cong
\mathrm{Im}\,\HH\otimes\CC$.  Gluons arise as the 8 generators of $\su(3)$.
\end{theorem}

\begin{proof}
$\mathrm{Im}\,\HH = \mathrm{span}_\RR\{\mathbf{e}_1,\mathbf{e}_2,\mathbf{e}_3\} \cong \RR^3$
under the inner product $\langle\mathbf{e}_i,\mathbf{e}_j\rangle = \delta_{ij}$.
The orientation-preserving isometries of $\RR^3$ form $\mathrm{SO}(3) \cong
\mathrm{SU}(2)/\ZZ_2$.  Complexifying: $\mathrm{Im}\,\HH\otimes_\RR\CC \cong \CC^3$
and the isometries extend to $\mathrm{SU}(3)$ (the compact automorphism group
of the Hermitian inner product on $\CC^3$).
The commutation with all three involutions is the constraint that selects
$\mathrm{SU}(3)$ from the larger $\mathrm{U}(3)$.
See \texttt{canonical/su3\_derivation/} and \texttt{canonical/interactions/sm\_gauge.tex §G.B--G.D}.
\end{proof}

\subsection{Quark representations and structural confinement}

\begin{corollary}[Quark multiplet, \PROVED\ \Lzero]
\label{cor:quarks}
Fields in the fundamental $\mathbf{3}$ representation of SU(3) are the
quarks.  Fields in the adjoint $\mathbf{8}$ representation are gluons.
Anti-quarks are in the conjugate $\bar{\mathbf{3}}$.
The confinement of colour-charged states to colour-singlet hadrons is a
structural property: the UBT action $S[\Theta]$ on the $\mathrm{SU}(3)_c$-invariant
sector produces only colour-singlet observables (structural confinement).
\end{corollary}

\begin{remark}[Dynamical confinement]
Structural confinement (colour-singlet observables) follows from the
SU(3) representation structure.  \emph{Dynamical} confinement (the linear
rising potential at large distances, the mass gap) is a separate problem
and corresponds to the Clay Millennium Prize problem; it is not claimed here.
\end{remark}

% ==============================================================================
\section{SU(2)$_L$ and Chirality}
\label{sec:su2}
% ==============================================================================

\subsection{SU(2) from left quaternion multiplication}

\begin{theorem}[SU(2)$_L$ from left action, \PROVED\ \Lzero]
\label{thm:su2}
The group $\mathrm{SU}(2)_L$ acts on $\BB = \CC\otimes\HH$ by left multiplication:
\[
  g\cdot(\CC\otimes q) \;=\; \CC\otimes(g\,q), \qquad g \in \mathrm{SU}(2) \subset \HH.
\]
The generators are $\tau^i = \sigma^i/2$ ($\sigma^i$ = Pauli matrices) acting
on the left-handed spinor representation $\mathbf{2}_L$.
The gauge bosons $W^\pm, W^3$ arise as the three generators of $\su(2)_L$.
\end{theorem}

\begin{proof}
$\HH \cong \CC^2$ as a left $\mathrm{SU}(2)$-module, with generators
$1\otimes\mathbf{e}_1, 1\otimes\mathbf{e}_2, 1\otimes\mathbf{e}_3$
represented by $-i\sigma_i$ in the $2\times 2$ complex matrix basis.
The isomorphism $\CC\otimes\HH \cong \mathrm{Mat}(2,\CC)$ maps left multiplication
by $\mathrm{SU}(2)$ to the standard $\mathrm{SU}(2)_L$ action on $\CC^2$.
See \texttt{canonical/interactions/sm\_gauge.tex §SU2}.
\end{proof}

\subsection{Chirality: Gap C1 closed}

\begin{theorem}[Chirality from $P_\psi$-involution, \PROVED\ \Lone]
\label{thm:chirality}
The involution $P_\psi: \psi \to -\psi$ on the imaginary-time circle $S^1_\psi$
acts on $\BB$ as a parity operator.  The $P_\psi$-odd sector of the UBT field
$\Theta$ defines left-chiral doublets that couple exclusively to $W^\pm$;
the $P_\psi$-even sector defines right-chiral singlets that do not couple to $W^\pm$.

This is the UBT mechanism for parity violation in weak interactions.
The absence of a right-handed $W_R$ coupling in the minimal action $S[\Theta]$
is a structural selection rule, not a free choice.
\end{theorem}

\begin{proof}
The UBT action $S[\Theta]$ contains the term
$\Tr[(D_\mu\Theta)^\dagger(D^\mu\Theta)]$ with $D_\mu$ including the
left $\mathrm{SU}(2)_L$ connection.  The $P_\psi$-involution anti-commutes
with the $\mathrm{SU}(2)_L$ generator acting on the left-chiral component.
No right-handed SU(2) coupling appears in the minimal action.
The chirality of the $W^\pm$ vertex follows from $P_\psi(\calL_W) \neq \calL_W$.
See \texttt{canonical/chirality/step3\_gap\_C1\_resolution.tex}.
The no-$W_R$ axiom is the minimal-coupling assumption for $S[\Theta]$.
\end{proof}

\begin{remark}[Gap C1 status]
Theorem~\ref{thm:chirality} closes \textbf{Gap~C1}: the $P_\psi$-involution
selects left-chiral coupling of the $W^\pm$ bosons.  This is the UBT
explanation of parity violation in weak interactions.  The proof is conditional
on the no-$W_R$ axiom; the deeper derivation of this axiom from $\CC\otimes\HH$
without invoking minimality is an open problem (item OP-S3 in \texttt{open\_problems.md}).
\end{remark}

% ==============================================================================
\section{U(1)$_Y$ Hypercharge and Electromagnetism}
\label{sec:u1}
% ==============================================================================

\subsection{U(1)$_Y$ from right phase action}

\begin{theorem}[U(1)$_Y$ from right phase, \PROVED\ \Lzero]
\label{thm:u1y}
The group $\mathrm{U}(1)_Y$ acts on $\CC\otimes\HH$ by right multiplication
by a complex phase:
\[
  e^{i\alpha}\cdot(\CC\otimes q) \;=\; \CC\otimes(q\,e^{i\alpha/2}).
\]
The generator is $Y = \frac{1}{2}\mathbf{I}$ (the unit $2\times 2$ matrix times $1/2$),
acting as the hypercharge operator on $\mathrm{Mat}(2,\CC) \cong \CC\otimes\HH$.
\end{theorem}

\begin{proof}
The right $\mathrm{U}(1)$ action commutes with the left $\mathrm{SU}(2)_L$ action
(by definition of left vs right multiplication) and with $\mathrm{SU}(3)_c$
(which acts on $\mathrm{Im}\,\HH$, not on the phase).  The generator $Y = \frac{1}{2}\mathbf{I}$
in the $\mathrm{Mat}(2,\CC)$ representation follows from the normalisation
convention consistent with the Gell-Mann--Nishijima relation.
See \texttt{canonical/interactions/sm\_gauge.tex §U1}.
\end{proof}

\subsection{Gell-Mann--Nishijima relation}

\begin{theorem}[Gell-Mann--Nishijima, \PROVED\ \Lone]
\label{thm:gmn}
In the UBT framework, the electromagnetic charge operator is
\[
  Q = T_3 + \frac{Y}{2},
\]
where $T_3 = \tau^3/2 = \sigma^3/4$ is the third generator of $\mathrm{SU}(2)_L$
and $Y$ is the hypercharge operator.  After electroweak symmetry breaking,
$\mathrm{U}(1)_{\mathrm{EM}}$ corresponds to the diagonal U(1) generated by $Q$.
\end{theorem}

\begin{proof}
The photon field $A_\mu = \sin\theta_W W^3_\mu + \cos\theta_W B_\mu$ is the
linear combination of $\mathrm{SU}(2)_L$ and $\mathrm{U}(1)_Y$ gauge fields
that remains massless after electroweak symmetry breaking.  The coupling of
$A_\mu$ to fermions defines $Q$.  In the standard electroweak basis this gives
the Gell-Mann--Nishijima relation.  See \texttt{canonical/qed.tex} and
\texttt{canonical/interactions/sm\_gauge.tex §U1}.
\end{proof}

\subsection{Hypercharge quantisation (partial result on Gap~C2)}

\begin{theorem}[Hypercharge quantisation from Dirac condition, \PROVED\ \Lzero\ + \Lone]
\label{thm:hcq}
The Dirac quantisation condition on $S^1_\psi$ forces
\[
  Y \;\in\; \tfrac{1}{6}\ZZ \quad\text{(for colour-charged fields)}, \qquad
  Y \;\in\; \tfrac{1}{2}\ZZ \quad\text{(for colour-neutral fields)}.
\]
The standard SM assignments $Y_Q = 1/6$, $Y_u = 2/3$, $Y_d = -1/3$,
$Y_L = -1/2$, $Y_e = -1$, $Y_\nu = 0$ are all consistent with these constraints.
\end{theorem}

\begin{proof}
Theorem~3.1 of \texttt{canonical/interactions/hypercharge\_assignments.tex}:
the Dirac condition gives $Y \in \frac{1}{2}\ZZ$ for colour-neutral fields.
For colour-charged fields (quarks in $\mathbf{3}$), the SU(3) lattice introduces
a factor of $1/3$, giving $Y \in \frac{1}{6}\ZZ$.
All standard SM hypercharge values are integer multiples of $1/6$, hence consistent.
\end{proof}

\begin{remark}[Gap C2 status]
Theorem~\ref{thm:hcq} shows that $Y_Q = 1/6$ is \emph{consistent} with the
UBT algebraic structure.  The stronger claim — that $Y_Q = 1/6$ is the
\emph{unique} admissible value forced by $\CC\otimes\HH$ without reference
to experimental charge assignments — is Gap~C2, currently at level \MC.
See \texttt{canonical/interactions/hypercharge\_assignments.tex} for the
partial derivation and honest verdict.
\end{remark}

% ==============================================================================
\section{Three Generations}
\label{sec:generations}
% ==============================================================================

\begin{theorem}[Three generations, \PROVED\ \Lzero]
\label{thm:generations}
The number of fermion generations is
\[
  N_{\mathrm{gen}} \;=\; \dim_\RR(\mathrm{Im}\,\HH) \;=\; 3.
\]
The three winding modes on $S^1_\psi$ carry identical $\mathrm{SU}(3)_c\times
\mathrm{SU}(2)_L\times\mathrm{U}(1)_Y$ quantum numbers and correspond to the
three known quark and lepton generations.
\end{theorem}

\begin{proof}
$\mathrm{Im}\,\HH = \mathrm{span}_\RR\{\mathbf{e}_1,\mathbf{e}_2,\mathbf{e}_3\}$
has real dimension~3 (\Lzero\ from definition of $\HH$).
Each imaginary direction supports a distinct winding sector on $S^1_\psi$;
these three sectors are related by the $\mathrm{SU}(3)$ colour action
(Theorem~\ref{thm:su3}) and therefore carry identical gauge quantum numbers.
The identification $N_{\mathrm{gen}} = 3$ follows.
See \texttt{DERIVATION\_INDEX.md §ψ-modes} and \texttt{canonical/su3\_derivation/}.
\end{proof}

\begin{remark}
The mass hierarchy between the three generations ($m_e \ll m_\mu \ll m_\tau$)
is not derived from $N_{\mathrm{gen}} = 3$.  The KK-mismatch theorem
(proved in \texttt{ARCHIVE/archive\_legacy/deprecated/research\_tracks/fermion\_mass\_program.md})
shows that the torus approach to mass ratios encounters an obstruction.
Generation masses are open problems (§\ref{sec:open}).
\end{remark}

% ==============================================================================
\section{Honest Dead Ends}
\label{sec:deadends}
% ==============================================================================

\subsection{Weinberg angle: declared dead end}

\begin{tcolorbox}[enhanced,colback=red!4!white,colframe=red!60!black,
  arc=3pt,boxrule=1.2pt,title={\textbf{Gap EW (Dead End): Weinberg Angle}}]

The following statement is copied verbatim from
\texttt{reports/gauge\_truth\_matrix.md §6}:

\medskip
\textbf{Gap EW (dead end)}: The biquaternion algebra $\CC\otimes\HH$
contains both $\mathrm{SU}(2)_L$ and $\mathrm{U}(1)_Y$ as subgroups, but
does not fix the ratio $g'/g = \tan\theta_W$ of their coupling constants.
The two couplings $g$ and $g'$ appear as independent parameters in any
Lagrangian on $\CC\otimes\HH$.  No purely algebraic argument has been found
to enforce $\sin^2\theta_W = 0.231$, and a no-go argument suggests none
exists: the algebra admits continuous deformations of the $\mathrm{SU}(2)_L
\times\mathrm{U}(1)_Y$ embedding that change $\tan\theta_W$ continuously.
The Weinberg angle is therefore a semi-empirical input in UBT at this stage.
\end{tcolorbox}

\noindent
This paper does \emph{not} claim the Weinberg angle is derived.  Any future
version of UBT that claims to derive $\sin^2\theta_W$ must overcome the
no-go argument stated above.

\medskip
\textbf{Update (2026-05-11)}: The EW1+RG route via generator norm ratios
gives $\sin^2\theta_W^{\rm GUT} = 3/8$ conditionally on the derivation of
fermion hypercharge assignments from $\CC\otimes\HH$ (Gap~C2).
The claim is upgraded from \emph{dead end} to \emph{conditional [L1]}
pending the colour-charge lattice theorem;
see \texttt{research\_tracks/EW/weinberg\_angle\_ew1\_rg.tex}.

\subsection{Dynamical confinement: not claimed}

The Clay Millennium problem of proving that SU(3) gauge theory has a mass gap
and confines colour is not addressed in this paper.  Structural confinement
(colour-singlet observables) follows from the SU(3) representation structure
(Corollary~\ref{cor:quarks}); dynamical confinement does not.

\subsection{$N_{\mathrm{eff}} = 12$: under audit in the companion alpha track}

The structural count $N_{\mathrm{eff}} = \dim G_{\mathrm{SM}} = 8+3+1 = 12$
(see \texttt{research\_tracks/T3\_ALPHA/neff12\_derivations.tex}, Route~R2) is
proved at \Lzero\ as an algebraic identity.  However, the identification of
this count with the number of charged scalar fields entering the one-loop
$U(1)_{\mathrm{EM}}$ photon vacuum polarisation is \emph{not} proved in the
current canonical files and is under critical audit in the companion T3\_ALPHA
track (\texttt{canonical/n\_eff/neff\_reconciliation.tex}).

The present gauge paper does not use $N_{\mathrm{eff}} = 12$ as an input
to $\alpha$ or any quantitative prediction.  It appears only as a structural
observation ($\dim G_{\mathrm{SM}} = 12$).  If the T3\_ALPHA audit revises
this count, no result in the present paper is affected.

Concretely, the ongoing revision of $N_{\mathrm{eff}}(\mathrm{loop})$ in the
alpha programme does not modify any theorem in Sections~\ref{sec:algebra}--
\ref{sec:generations}: those statements are algebraic/representation-theoretic
and are anchored in the gauge-generator count $N_{\mathrm{eff}}(\mathrm{R2})=12$.
Therefore the T2\_GAUGE conclusions remain unchanged by the loop-audit update.

For completeness, the effective coupling
$B\approx 46.28$ (Gap~G137-B) is not derived from $S[\Theta]$.
The numerical observation
$B=N_{\mathrm{eff}}^{3/2}\cdot 2^{1/8}\cdot\vartheta_3(0|i)^{1/4}$
has accuracy about $0.007\%$ but lacks a first-principles derivation.
A Seeley-DeWitt mechanism is under investigation.

% ==============================================================================
\section{Open Problems}
\label{sec:open}
% ==============================================================================

The following problems are identified as open in the UBT gauge sector.
No claim is made about any of them in this paper.

\begin{description}[style=nextline]
  \item[Gap~C2: Specific hypercharge assignments]
    Derive $Y_Q = 1/6$ (rather than showing consistency) from
    $\CC\otimes\HH$ without using experimental charge values as input.
    See \texttt{canonical/interactions/hypercharge\_assignments.tex} for the
    partial result (Dirac quantisation forces $Y \in \frac{1}{6}\ZZ$;
    uniqueness of $Y_Q = 1/6$ is \MC).

  \item[Gap~EW-2: Higgs doublet from $S[\Theta]$]
    The Higgs field and its symmetry-breaking VEV have not been derived
    from the UBT action.  The pattern $\mathrm{SU}(2)_L\times\mathrm{U}(1)_Y
    \to \mathrm{U}(1)_{\mathrm{EM}}$ is motivated but not derived.

  \item[Fermion mass spectrum]
    The formula $m(n) = A\cdot n^p - B_m\cdot n\ln n$ reproduces $m_e$ to
    $0.22\%$ accuracy (semi-empirical \textbf{[SE]}), but the parameter
    $B_m \approx -14.1\,\mathrm{MeV}$ is fitted, not derived.
    The ratios $m_\mu/m_e$, $m_\tau/m_\mu$ are unexplained.
    See \texttt{research\_tracks/fermion\_masses/fermion\_mass\_status\_v2.md}.

  \item[Yukawa couplings and CKM/PMNS]
    The Yukawa coupling matrix $y_{ij}$ and the CKM/PMNS mixing matrices
    are not derived in UBT.  These are deferred to a future programme.

  \item[Strong coupling $g_s$]
    The strong coupling constant $g_s$ (or equivalently $\alpha_s$) at any
    scale is not derived from $\CC\otimes\HH$ at this stage.
\end{description}

% ==============================================================================
\section{Summary}
\label{sec:summary}
% ==============================================================================

\begin{center}
\begin{tabular}{lll}
\toprule
\textbf{Result} & \textbf{Level} & \textbf{Status} \\
\midrule
$\CC\otimes\HH \cong \mathrm{Mat}(2,\CC)$ & \Lzero & \PROVED \\
$\mathrm{SU}(3)_c$ from involutions & \Lzero & \PROVED \\
$\mathrm{SU}(2)_L$ from left action & \Lzero & \PROVED \\
$\mathrm{U}(1)_Y$ from right phase & \Lzero & \PROVED \\
Gell-Mann--Nishijima $Q = T_3 + Y/2$ & \Lone & \PROVED \\
$N_{\mathrm{gen}} = 3$ & \Lzero & \PROVED \\
Chirality (Gap~C1 closed) & \Lone & \PROVED \\
$Y \in \frac{1}{6}\ZZ$ (Dirac quantisation) & \Lzero+\Lone & \PROVED \\
\midrule
Weinberg angle $\sin^2\theta_W$ & — & \DEADEND (stated explicitly) \\
Higgs doublet VEV & — & \OPEN\ (Gap EW-2) \\
Fermion masses & — & \OPEN\ (semi-empirical at best) \\
$Y_Q = 1/6$ unique & — & \MC\ (partial; Gap~C2 not closed) \\
\bottomrule
\end{tabular}
\end{center}

Together with the companion GR paper~\cite{jaros2026_gr}, UBT accounts for
the gauge structure $\mathrm{SU}(3)_c\times\mathrm{SU}(2)_L\times\mathrm{U}(1)_Y$
and General Relativity from the single algebraic object $\CC\otimes\HH$.
Significant open problems — particularly the Weinberg angle, fermion masses,
and the Higgs mechanism — remain and are stated explicitly.

% ==============================================================================
% References
% ==============================================================================

\begin{thebibliography}{99}

\bibitem{jaros2026_gr}
I.\ Jaro\v{s},
\emph{General Relativity as a Real-Projected Limit of Unified Biquaternion Theory},
preprint (2026), \texttt{papers/UBT\_GR\_Submission.tex}.

\bibitem{porteous1995}
I.\ R.\ Porteous, \emph{Clifford Algebras and the Classical Groups},
Cambridge University Press, 1995.

\bibitem{peskin_schroeder}
M.\ E.\ Peskin and D.\ V.\ Schroeder,
\emph{An Introduction to Quantum Field Theory},
Addison-Wesley, 1995.

\bibitem{connes2008}
A.\ Connes and M.\ Marcolli,
\emph{Noncommutative Geometry, Quantum Fields and Motives},
AMS, 2008.

\bibitem{baez2010}
J.\ Baez and J.\ Huerta,
\emph{The Algebra of Grand Unified Theories},
Bull.\ Amer.\ Math.\ Soc.\ \textbf{47}, 483--552 (2010).

\bibitem{zerilli1970}
F.\ J.\ Zerilli,
\emph{Effective Potential for Even-Parity Regge-Wheeler Gravitational Perturbation Equations},
Phys.\ Rev.\ Lett.\ \textbf{24}, 737 (1970).

\bibitem{regge_wheeler1957}
T.\ Regge and J.\ A.\ Wheeler,
\emph{Stability of a Schwarzschild Singularity},
Phys.\ Rev.\ \textbf{108}, 1063 (1957).

\bibitem{gell_mann_nishijima}
M.\ Gell-Mann, Nuovo Cim.\ \textbf{4}, 848 (1956);
K.\ Nishijima, Prog.\ Theor.\ Phys.\ \textbf{13}, 285 (1955).

\end{thebibliography}

\end{document}
